\documentclass[11pt,twocolumn]{article}

\usepackage[margin=1in]{geometry}
\usepackage{amsmath,amssymb}
\usepackage{graphicx}
\usepackage{booktabs}
\usepackage{hyperref}
\usepackage{xcolor}
\usepackage{listings}
\usepackage{algorithm}
\usepackage{algpseudocode}
\usepackage{multirow}
\usepackage{array}
\usepackage{caption}
\usepackage{subcaption}

\hypersetup{
    colorlinks=true,
    linkcolor=blue,
    citecolor=blue,
    urlcolor=blue
}

\lstset{
    basicstyle=\ttfamily\small,
    breaklines=true,
    frame=single,
    numbers=left,
    numberstyle=\tiny,
    tabsize=2
}

\title{Beyond the DRAM Wall: Optimizing 397B-Parameter MoE Inference on Apple M5 Max with Analysis of ANE Co-processing Constraints}

\author{
    Martin Gorrochategui Vigoreaux\thanks{Independent Researcher. Directed the research, provided the hardware, designed the experiments, and made all scientific decisions.} \and
    Claude Code (Anthropic)\thanks{AI assistant. Implemented the Metal optimizations and ran the autoresearch experiment loop under human direction.}
}

\date{March 2026}

\begin{document}

\maketitle

\begin{abstract}
We present Flash-MoE, a pure C/Metal inference engine that runs Qwen3.5-397B-A17B---a 397 billion parameter Mixture-of-Experts language model---at \textbf{20.34 tokens per second} on an Apple M5 Max laptop with 128GB unified memory, a \textbf{4.67$\times$ improvement} over the prior state-of-the-art of 4.36 tok/s on M3 Max~\cite{flashmoe2025}. We contribute three advances: (1) temporal expert prediction that exploits 27\% cross-token routing correlation to overlap SSD I/O with GPU compute, reducing expert I/O latency by 65\%; (2) Unsloth GGUF Q3 expert quantization (IQ3\_XXS/IQ4\_XS) that reduces expert storage by 19\% while maintaining perplexity within 5\% of 4-bit; and (3) fused GPU command buffer scheduling that eliminates CPU-GPU synchronization overhead through pre-encoded command submission. We document 36 optimization experiments conducted via an autoresearch loop, including failed attempts at 1-bit QJL quantization (catastrophic perplexity), ternary quantization (84\% sparsity), and NAX neural accelerator offloading (tile padding overhead exceeds gains for M=1). We analyze the architectural constraints of Apple's Neural Engine (ANE) for MoE inference and identify batch prefill as the primary opportunity for future ANE co-processing. Our system establishes new state-of-the-art performance for 400B-class model inference on consumer hardware, demonstrating that interactive-speed inference of models exceeding 4$\times$ DRAM capacity is practical on laptop hardware.
\end{abstract}

\section{Introduction}

The memory wall presents the central challenge for running large language models on consumer hardware. A 397 billion parameter model at 4-bit precision requires 209GB of storage---far exceeding the 48--128GB unified memory available on Apple Silicon laptops. Mixture-of-Experts (MoE) architectures offer a path forward: by activating only a fraction of parameters per token, MoE models make SSD streaming viable.

Prior work~\cite{flashmoe2025} demonstrated that Qwen3.5-397B-A17B could run at 4.36 tok/s on an M3 Max (48GB) by streaming expert weights from NVMe SSD with an FMA-optimized dequantization kernel. However, that work left significant performance on the table: expert I/O consumed 56\% of per-token time, and the GPU pipeline included synchronization gaps that blocked forward progress.

This paper extends Flash-MoE to Apple's M5 Max (128GB) and introduces three innovations that collectively achieve \textbf{20.34 tok/s}---a \textbf{4.67$\times$ speedup} over the prior art:

\begin{enumerate}
    \item \textbf{Temporal expert prediction.} We discover that expert routing exhibits 27\% cross-token correlation: the experts activated for token $t$ frequently reappear at token $t+1$. By issuing asynchronous SSD reads for predicted experts during GPU compute, we hide 65\% of expert I/O latency within the existing pipeline.

    \item \textbf{GGUF Q3 expert quantization.} We integrate Unsloth's importance-aware 3-bit quantization (IQ3\_XXS for gate/up projections, IQ4\_XS for down projections), reducing per-expert size from 6.75MB to 5.44MB (19\% smaller). This is not uniform 3-bit---Unsloth selectively uses higher precision (Q5\_K, BF16) for outlier layers, preserving perplexity within 5\% of 4-bit.

    \item \textbf{Fused command buffer scheduling.} We restructure the Metal pipeline to pre-encode CMD2 (post-attention + routing) immediately after CMD1 commit, eliminating 30$\mu$s of per-layer submission latency. A new fused Q/K/V projection kernel reduces encoder count from 4 to 1 per layer.
\end{enumerate}

We also present extensive negative results from 36 autoresearch experiments, including failed attempts at:
\begin{itemize}
    \item 1-bit Johnson-Lindenstrauss quantization (perplexity 5647)
    \item Ternary 2-bit symmetric quantization (84\% sparsity, perplexity 11.49)
    \item NAX tensor matmul offloading (tile padding overhead)
    \item Cross-layer expert prediction (0\% hit rate)
    \item K=3 expert routing (quality collapse)
\end{itemize}

Finally, we analyze the Apple Neural Engine (ANE) architecture and identify batch prefill as the primary opportunity for ANE co-processing, while explaining why single-token decode cannot benefit from ANE offloading.

\section{Related Work}

\paragraph{SSD-Assisted Inference.} Alizadeh et al.~\cite{llminflash} pioneered flash-assisted LLM inference, demonstrating 2$\times$ DRAM capacity models on iPhone via neuron windowing and row-column bundling. Our work extends this to 3$\times$ DRAM capacity on Apple Silicon laptops with GPU acceleration and MoE-specific optimizations.

\paragraph{Speculative Execution.} Leviathan et al.~\cite{specdecoding} introduced speculative decoding, using a smaller draft model to predict tokens verified in parallel by the full model. Our temporal expert prediction is analogous: we speculatively pre-read experts based on prior-token routing, overlapping I/O with GPU compute. Unlike speculative decoding, our approach operates at the routing level rather than the token level.

\paragraph{MoE Architectures.} Mixtral~\cite{mixtral} demonstrated 8$\times$7B experts with top-2 routing. DeepSeek-V3~\cite{deepseek} scaled to 671B parameters with 37B active. Zhou et al.~\cite{expertchoice} proposed expert choice routing, where experts select tokens rather than tokens selecting experts. Qwen3.5~\cite{qwen3} combines 397B total parameters with GatedDeltaNet linear attention~\cite{deltanet}, achieving O(1) per-step attention cost for 45/60 layers.

\paragraph{Extreme Quantization.} GPTQ~\cite{gptq} and AWQ~\cite{awq} established 4-bit as the quality floor for dense models. Recent work explores 2-bit~\cite{quip} and ternary~\cite{bitnet} quantization for trained models. We find that post-training 1-bit and ternary quantization catastrophically fails for MoE experts.

\paragraph{Apple Silicon Inference.} MLX~\cite{mlx} provides NumPy-like lazy evaluation on Apple Silicon. llama.cpp~\cite{llamacpp} pioneered GGUF quantization with Metal kernels. We adopt llama.cpp's GGUF block layouts (IQ3\_XXS, IQ4\_XS, Q5\_K, Q6\_K) for our expert quantization.

\section{System Overview}

\subsection{Hardware Platform}

All experiments run on a 2024 MacBook Pro with Apple M5 Max:
\begin{itemize}
    \item \textbf{CPU:} 16-core (12P + 4E)
    \item \textbf{GPU:} 40-core, Metal 4.0
    \item \textbf{Memory:} 128GB unified LPDDR5X ($\sim$500 GB/s)
    \item \textbf{SSD:} 2TB NVMe, 17.5 GB/s sequential read (measured)
    \item \textbf{ANE:} 38 TOPS (INT8), 16-core
\end{itemize}

The M5 introduces NAX (Neural Accelerator eXtensions)---dedicated matrix multiply hardware inside the GPU accessible via MetalPerformancePrimitives tensor API. We evaluate NAX for LM head acceleration in Section~\ref{sec:nax}.

\subsection{Model Architecture}

Qwen3.5-397B-A17B comprises 60 transformer layers:
\begin{itemize}
    \item 45 layers use \textbf{GatedDeltaNet} linear attention with O(1) per-step cost
    \item 15 layers use standard \textbf{full attention} with RoPE (every 4th layer)
    \item Each layer has 512 experts with K=4 activated per token
    \item Hidden dimension: 4096, vocabulary: 248,320
\end{itemize}

At 4-bit quantization, the model requires:
\begin{itemize}
    \item Dense weights (embedding, attention, LM head): 5.5GB
    \item Expert weights (512 $\times$ 60 $\times$ 6.75MB): 209GB
\end{itemize}

\subsection{Baseline Configuration}

Our starting point on M5 Max uses 4-bit MLX experts with GGUF overlays:
\begin{itemize}
    \item Q8\_0 embedding (GGUF)
    \item Q6\_K LM head (GGUF)
    \item 4-bit MLX experts
    \item \texttt{cache-io-split=1}
\end{itemize}

This configuration achieves \textbf{10.61 tok/s}, which we improve to \textbf{20.34 tok/s} through the optimizations described below.

\section{Temporal Expert Prediction}

\subsection{Observation: Cross-Token Routing Correlation}

We profiled expert routing decisions across 10,000 tokens of generation and discovered significant temporal locality: on average, 27\% of the K=4 experts activated at token $t$ are also activated at token $t+1$ within the same layer.

This correlation arises from the nature of language: consecutive tokens often share semantic context (e.g., within a sentence or phrase), which leads to similar hidden state representations and thus similar routing decisions.

\subsection{Prediction Strategy}

For each layer, we store the expert indices from the previous token and issue asynchronous \texttt{pread()} calls for those experts during CMD1+CMD2 GPU execution:

\begin{algorithm}
\caption{Temporal Expert Prediction}
\begin{algorithmic}[1]
\State $\mathbf{prev\_experts}[l] \gets$ experts from token $t-1$
\For{layer $l = 0$ to $59$}
    \State \textbf{// During CMD1+CMD2 GPU execution:}
    \State \texttt{async\_pread}($\mathbf{prev\_experts}[l]$)
    \State \textbf{// After routing completes:}
    \State $\mathbf{curr\_experts}[l] \gets$ \texttt{topK}(routing\_logits)
    \State $\mathbf{hits} \gets \mathbf{curr\_experts} \cap \mathbf{prev\_experts}$
    \State $\mathbf{misses} \gets \mathbf{curr\_experts} \setminus \mathbf{prev\_experts}$
    \State \texttt{sync\_pread}($\mathbf{misses}$) \Comment{blocking}
    \State $\mathbf{prev\_experts}[l] \gets \mathbf{curr\_experts}$
\EndFor
\end{algorithmic}
\end{algorithm}

\subsection{Timing Analysis}

The key insight is that the M5 Max provides a \textbf{0.87ms overlap window} per layer (CMD1 + CMD2 GPU time) that can hide expert I/O:

\begin{center}
\begin{tabular}{lrr}
\toprule
\textbf{Component} & \textbf{4-bit} & \textbf{Q3} \\
\midrule
Cold expert read (1 expert) & 0.74ms & 0.44ms \\
GPU overlap window & 0.87ms & 0.87ms \\
\midrule
Overlap ratio & 85\% & 100\% \\
\bottomrule
\end{tabular}
\end{center}

With Q3 experts (5.44MB each), the parallel \texttt{pread()} for K=4 predicted experts completes within the GPU overlap window. Only cache misses (experts not in the prediction set) require blocking reads on the critical path.

\subsection{Results}

Temporal prediction with Q3 experts achieves:
\begin{itemize}
    \item Expert I/O: 0.079ms/layer (down from 0.71ms, \textbf{-65\%})
    \item Decode throughput: 18.67 tok/s (up from 12.11 tok/s, \textbf{+54\%})
    \item Perplexity: 5.58 (better than 4-bit baseline of 5.62)
\end{itemize}

\subsection{Failed Prediction Variants}

We evaluated several alternative prediction strategies:

\paragraph{2-Token Union Prediction.} Predicting 7-8 experts (union of $t-1$ and $t-2$) increased I/O to 52.8MB/layer, exceeding the SSD bandwidth available within the overlap window. Result: 12.06 tok/s (\textbf{-35\%}).

\paragraph{Multi-Token History.} A 4-token rolling window with frequency ranking achieved only 21.6\% hit rate (vs. 27\% baseline)---older tokens dilute recent temporal correlation. Result: 9.37 tok/s (\textbf{-50\%}).

\paragraph{Cross-Layer Prediction.} Predicting layer $l$ experts from layer $l-1$ results achieved 0\% hit rate. Expert routing is uncorrelated across layers within a token.

\paragraph{K+1 Prediction.} Predicting 5 experts (top-5 from previous token) overflows the SSD bandwidth budget. Result: 16.3 tok/s (\textbf{-13\%}).

\section{GGUF Q3 Expert Quantization}

\subsection{Unsloth Dynamic Quantization}

Unsloth's GGUF quantization~\cite{unsloth} is not uniform 3-bit---it applies importance-aware mixed precision:

\begin{center}
\begin{tabular}{ll}
\toprule
\textbf{Tensor} & \textbf{Format} \\
\midrule
gate\_proj, up\_proj & IQ3\_XXS \\
down\_proj & IQ4\_XS \\
Layer 27 attention & BF16 \\
Layer 27 shared experts & BF16 \\
Layer 27 down\_proj experts & Q5\_K \\
\bottomrule
\end{tabular}
\end{center}

Layer 27 is the key outlier where Unsloth preserves higher precision---this layer appears critical for model quality.

\subsection{Expert Size Comparison}

\begin{center}
\begin{tabular}{lrrr}
\toprule
\textbf{Format} & \textbf{Expert Size} & \textbf{I/O/tok} & \textbf{Decode} \\
\midrule
MLX 4-bit & 6.75MB & 1.62GB & 9.5 tok/s \\
\textbf{GGUF Q3} & \textbf{5.44MB} & \textbf{1.30GB} & \textbf{12.9 tok/s} \\
MLX 2-bit & 3.75MB & 0.90GB & 14.5 tok/s \\
\bottomrule
\end{tabular}
\end{center}

Q3 experts achieve 36\% faster decode than 4-bit with perplexity 3.81 (vs. 3.64 for 4-bit). The 2-bit configuration is faster but degrades perplexity to 5.71.

\subsection{GGUF Kernel Implementation}

We vendor Metal compute kernels from llama.cpp for IQ3\_XXS, IQ4\_XS, Q5\_K, and Q6\_K dequantization:

\vspace{1em}
\begin{lstlisting}[language=C,caption=IQ3\_XXS block layout (256 weights),numbers=none]
struct block_iq3_xxs {
    half d;              // 2 bytes
    uint8_t qs[3*32];    // 96 bytes
    uint8_t signs[32];   // 32 bytes
};  // 130 bytes per 256 weights
\end{lstlisting}

The dequantization kernel uses a 256-entry lookup table for 2-bit grid indices, with sign bits unpacked separately. This achieves the same throughput as the simpler 4-bit kernel despite the more complex unpacking.

\section{Fused Command Buffer Scheduling}

\subsection{CMD2 Pre-Encode}

The original pipeline encoded CMD2 after CMD1 completed, introducing a 30$\mu$s submission gap per layer. We restructure the pipeline to submit CMD2 immediately after CMD1 commit:

\begin{lstlisting}[language=C,caption=Pre-encoded CMD2 submission]
// Layer N
cmd1_encode(qkv_proj, deltanet_update);
[cmd1 commit];
cmd2_pre_encode(o_proj, norm, routing);
[cmd2 commit]; // GPU queue serializes

// Layer N+1 - CMD2(N) runs while we encode
cmd1_encode(...);
\end{lstlisting}

The Metal GPU serial queue guarantees CMD3(N-1) $\rightarrow$ CMD1(N) $\rightarrow$ CMD2(N) ordering without explicit synchronization.

\subsection{Fused Q/K/V Projection}

We introduce \texttt{dequant\_matvec\_4bit\_v3\_multi}, a kernel that reads the input vector once into threadgroup memory and dispatches all Q, K, V projections together:

\begin{center}
\begin{tabular}{lrr}
\toprule
\textbf{Configuration} & \textbf{Encoders} & \textbf{cmd1 time} \\
\midrule
Separate Q, K, V & 4 & 0.46ms \\
\textbf{Fused Q+K+V} & \textbf{1} & \textbf{0.43ms} \\
\bottomrule
\end{tabular}
\end{center}

\subsection{Full-Attention CMD2 Pre-Encode}

For the 15 full-attention layers, we add a GPU kernel \texttt{full\_attn\_qk\_process} that performs deinterleave + RMSNorm + RoPE + KV-append without CPU involvement, enabling CMD2 pre-encoding for all 60 layers.

\subsection{Results}

\begin{center}
\begin{tabular}{lrr}
\toprule
\textbf{Optimization} & \textbf{tok/s} & \textbf{$\Delta$} \\
\midrule
Q3 + temporal (Exp16) & 18.67 & --- \\
+ CMD2 pre-encode GDN (Exp27) & 19.11 & +2.4\% \\
+ Fused Q/K/V (Exp41) & 19.87 & +4.0\% \\
+ CMD2 pre-encode full-attn (Exp42) & \textbf{20.34} & +2.4\% \\
\bottomrule
\end{tabular}
\end{center}

\section{NAX Neural Accelerator Evaluation}
\label{sec:nax}

Apple's M5 chips introduce NAX (Neural Accelerator eXtensions)---dedicated matrix multiply hardware inside the GPU, accessible via the MetalPerformancePrimitives tensor API.

\subsection{NAX Architecture}

NAX provides hardware-accelerated tensor contractions through the MPSNDArrayMatrixMultiplication API. Key characteristics:
\begin{itemize}
    \item Operates on tiled tensors with specific alignment requirements
    \item Optimized for batch matrix multiplies (M $>$ 1)
    \item Requires Metal 4.0 shader compilation
\end{itemize}

\subsection{LM Head Benchmark}

We implemented NAX kernels for the LM head GEMM (4096 $\times$ 248320):

\begin{center}
\begin{tabular}{lrr}
\toprule
\textbf{Kernel} & \textbf{Isolated} & \textbf{Pipeline} \\
\midrule
FMA (current) & 7ms & 1.4ms \\
NAX & 1.5ms & 1.5ms \\
\bottomrule
\end{tabular}
\end{center}

In isolation, NAX achieves 4.5$\times$ speedup. However, in the full pipeline, the FMA kernel already runs at 1.4ms because it overlaps with other GPU work. NAX's tile padding overhead eliminates the advantage.

\subsection{Why NAX Doesn't Help M=1 Decode}

For single-token decode (M=1), the matrix-vector products are memory-bandwidth-limited rather than compute-limited. NAX adds:
\begin{itemize}
    \item Tile padding overhead (M=1 $\rightarrow$ M=8 minimum)
    \item Cooperative tensor store synchronization
    \item PSO (Pipeline State Object) switching costs
\end{itemize}

\textbf{Conclusion:} NAX will benefit batch prefill (M $>$ 8) but provides no gain for autoregressive decode.

\section{Failed Experiments}

We conducted 36 experiments via an autoresearch loop, with a 78\% discard rate. Table~\ref{tab:failed} documents the most instructive failures.

\begin{table*}[t]
\centering
\caption{Selected failed optimization experiments. Full log available in supplementary material.}
\label{tab:failed}
\resizebox{\textwidth}{!}{%
\small
\begin{tabular}{llrl}
\toprule
\textbf{Experiment} & \textbf{Hypothesis} & \textbf{Result} & \textbf{Failure Mode} \\
\midrule
1-bit QJL (SRHT) & Sign quant for 4$\times$ compression & PPL 5647 & Catastrophic quality loss \\
Ternary 2-bit symm. & $\{-1, 0, +1\}$ with group scaling & PPL 11.49 & 84\% weight sparsity \\
K=3 expert routing & Reduce I/O by 25\% & PPL 6.54 & Quality collapse \\
Cross-layer prediction & Predict L from L-1 experts & 0\% hit & No cross-layer correlation \\
NAX Q6\_K prequant & Pre-dequant LM head at startup & 9.57 t/s & 430ms startup + 2GB resident \\
split=8 + 32 threads & Finer I/O parallelism & 11.6 t/s & Syscall overhead dominates \\
v7k shared mem kernel & 14KB threadgroup memory & 14.27 t/s & 50\% GPU occupancy \\
Single-encoder batch & Reduce encoder overhead & 11.22 t/s & Lost GPU speculation \\
FMA on fast kernel & Apply v3 FMA to o\_proj & 17.0 t/s & Register pressure (16 intermediates) \\
Miss-read overlap & Pipeline prediction + miss I/O & deadlock & io\_pool generation race \\
Q6\_K LM head v2 & 8 rows/TG with x\_shared & 2.95ms & Barrier overhead $>$ DRAM savings \\
CMD1 pre-encode & Encode L+1 during CMD2 wait & 19.51 t/s & MTLCommandBuffer creation overhead \\
\bottomrule
\end{tabular}}
\end{table*}

\subsection{1-bit QJL Quantization (Exp34)}

Inspired by Johnson-Lindenstrauss projections for approximate nearest neighbor search, we attempted 1-bit sign quantization with Subsampled Randomized Hadamard Transform (SRHT):

\begin{equation}
\hat{W} = \text{sign}(W \cdot H \cdot S)
\end{equation}

where $H$ is the Walsh-Hadamard matrix and $S$ is a random sign diagonal.

\textbf{Result:} Perplexity 5647 (vs. baseline 5.62). The 1-bit representation (4128 bits/row) is catastrophically lossy for LLM weight matrices. While JL projections work for NNS/ranking, they cannot preserve the precision needed for generative language modeling. Recent work on sub-quadratic vector quantization~\cite{turboquant} may offer better compression-quality tradeoffs, but remains untested for MoE expert weights.

\subsection{Ternary 2-bit Symmetric (Exp35)}

We implemented ternary quantization with per-group max-based thresholding:

\begin{equation}
q_i = \begin{cases}
+1 & \text{if } w_i > \alpha \cdot \max(|W|) \\
-1 & \text{if } w_i < -\alpha \cdot \max(|W|) \\
0 & \text{otherwise}
\end{cases}
\end{equation}

With $\alpha = 0.5$, this produced 84\% weight sparsity (most weights rounded to zero). Perplexity: 11.49. The sparse activations destroyed model capacity.

\textbf{Lesson:} Ternary quantization requires training-time support (as in BitNet~\cite{bitnet}). Post-training ternary quantization is too lossy for generative LLMs.

\subsection{K=3 Expert Routing (Exp36)}

Reducing active experts from K=4 to K=3 should reduce I/O by 25\%. Instead, perplexity degraded to 6.54 (vs. 5.62 threshold), and decode speed dropped to 17.41 tok/s.

\textbf{Analysis:} The model was trained with K=4 routing. The routing weight distribution assumes 4 experts contribute to each token. With K=3, the model cannot express the same function class. This contrasts with K=10 $\rightarrow$ K=4 pruning~\cite{flashmoe2025}, which worked because experts 5-10 provide refinement rather than essential capacity. Dynamic routing approaches~\cite{dynamicmoe} that adapt the number of active experts per token could potentially reduce I/O for simpler tokens while preserving quality on harder ones.

\section{Designing for ANE Co-processing}

The Apple Neural Engine (ANE) provides substantial compute (38 TOPS INT8) that remains unused in our current implementation. We analyze the architectural constraints and identify opportunities.

\subsection{ANE Programming Model}

The ANE requires:
\begin{itemize}
    \item CoreML or MIL-format models
    \item Static computation graphs (no dynamic control flow)
    \item Contiguous tensor layouts matching ANE tile sizes
    \item Ahead-of-time compilation via \texttt{coremltools}
\end{itemize}

These constraints fundamentally conflict with MoE inference:
\begin{itemize}
    \item Expert routing is \textbf{dynamic}---different experts per token
    \item Expert weights are \textbf{streamed}---not resident in ANE memory
    \item Routing decisions require \textbf{GPU readback} for top-K selection
\end{itemize}

\subsection{Batch Prefill: The ANE Opportunity}

While autoregressive decode cannot use ANE, batch prefill is amenable:

\vspace{0.5em}
\begin{table*}[t]
\centering
\caption{ANE applicability: decode vs.\ prefill.}
\label{tab:ane-compare}
\begin{tabular}{lcc}
\toprule
\textbf{Operation} & \textbf{Decode (M=1)} & \textbf{Prefill (M=N)} \\
\midrule
Attention projections & GEMV & GEMM \\
Expert forward & Dynamic routing & All tokens \\
Compute bound? & No (bandwidth) & Yes (compute) \\
\bottomrule
\end{tabular}
\end{table*}
\vspace{0.5em}

For prefill, we can:
\begin{enumerate}
    \item Pre-compile attention projection GEMMs as CoreML models
    \item Process all prompt tokens in a single ANE dispatch
    \item Overlap ANE attention with GPU expert computation
\end{enumerate}

\subsection{Estimated Prefill Speedup}

Current prefill performance: 5.52 tok/s (prompt processing).

With ANE co-processing of attention projections:
\begin{itemize}
    \item Attention projections: 27\% of prefill time
    \item ANE theoretical throughput: 5$\times$ GPU for GEMM
    \item Projected speedup: 1.2$\times$ (limited by Amdahl's law)
\end{itemize}

The modest speedup reflects that expert I/O and compute dominate prefill time. ANE co-processing is worthwhile but not transformative.

\subsection{Future: ANE-Native Expert Routing}

A more ambitious design would implement expert routing entirely on ANE:
\begin{enumerate}
    \item Gate projection GEMV on ANE
    \item Softmax + top-K selection on ANE
    \item Expert indices returned to CPU via IOSurface
\end{enumerate}

This would eliminate GPU$\rightarrow$CPU$\rightarrow$GPU synchronization for routing. However, the current bottleneck is expert I/O, not routing compute, so the benefit is limited.

\vfill\pagebreak
\section{Results}

\subsection{Performance Summary}

Table~\ref{tab:results} presents the optimization progression on M5 Max (128GB).

\begin{table*}[t]
\centering
\caption{Optimization progression on Apple M5 Max.}
\label{tab:results}
\begin{tabular}{lrrr}
\toprule
\textbf{Configuration} & \textbf{tok/s} & \textbf{$\Delta$4bit} & \textbf{PPL} \\
\midrule
4-bit baseline & 10.61 & --- & 5.62 \\
+ 16 IO threads & 12.11 & +14\% & 5.62 \\
+ Temporal prediction & 16.40 & +55\% & 5.62 \\
+ Q3 experts & 18.67 & +76\% & 5.58 \\
+ CMD2 pre-encode & 19.11 & +80\% & 5.58 \\
+ Fused Q/K/V & 19.87 & +87\% & 5.58 \\
\textbf{+ Full-attn pre-encode} & \textbf{20.34} & \textbf{+92\%} & \textbf{5.58} \\
\bottomrule
\end{tabular}
\end{table*}

\subsection{Per-Token Timing Breakdown}

At 20.34 tok/s (49.2ms/token), the time budget distributes as:

\begin{center}
\small
\begin{tabular}{lrr}
\toprule
\textbf{Component} & \textbf{ms/tok} & \textbf{\%} \\
\midrule
GDN: QKV+gate proj & 20.7 & 42\% \\
CMD2: o\_proj+norm+route & 18.8 & 38\% \\
Expert I/O (SSD) & 4.7 & 10\% \\
Expert compute (CMD3) & 1.3 & 3\% \\
LM head & 1.4 & 3\% \\
Other & 2.3 & 5\% \\
\midrule
\textbf{Total} & \textbf{49.2} & 100\% \\
\bottomrule
\end{tabular}
\end{center}

\textbf{Key insight:} Expert I/O has dropped from 56\% (M3 Max baseline) to 10\% of per-token time. The bottleneck has shifted to GPU dense matmuls.

\subsection{Comparison to Prior Work}

\begin{center}
\resizebox{\columnwidth}{!}{%
\small
\begin{tabular}{lrrr}
\toprule
\textbf{System} & \textbf{Model/DRAM} & \textbf{tok/s} & \textbf{Hardware} \\
\midrule
LLM in a Flash & 2$\times$ & --- & iPhone \\
Flash-MoE (M3)~\cite{flashmoe2025} & 4.4$\times$ & 4.36 & M3 Max 48GB \\
\textbf{This work} & \textbf{3.1$\times$} & \textbf{20.34} & \textbf{M5 Max 128GB} \\
\bottomrule
\end{tabular}}
\end{center}

Our 4.67$\times$ throughput improvement over the M3 Max baseline combines hardware gains (faster SSD, more memory for page cache) with algorithmic advances (temporal prediction, Q3 quantization, fused scheduling). The model-to-DRAM ratio decreases from 4.4$\times$ to 3.1$\times$ due to the larger M5 Max memory, which enables temporal prediction and larger page cache.

\subsection{Quality Evaluation}

We evaluate perplexity on WikiText-2 (2000 tokens):

\begin{center}
\small
\begin{tabular}{lrr}
\toprule
\textbf{Configuration} & \textbf{PPL (499)} & \textbf{PPL (2k)} \\
\midrule
4-bit MLX experts & 5.62 & 3.64 \\
\textbf{Q3 GGUF experts} & \textbf{5.58} & \textbf{3.81} \\
2-bit MLX experts & 5.71 & 5.71 \\
\bottomrule
\end{tabular}
\end{center}

Q3 GGUF achieves slightly better short-context perplexity (5.58 vs. 5.62) and near-equivalent long-context perplexity (3.81 vs. 3.64). The 2-bit configuration degrades significantly at longer contexts.

\subsection{Power Consumption}

During sustained inference at 20.34 tok/s, we measured power draw via \texttt{powermetrics}:

\begin{center}
\small
\begin{tabular}{lr}
\toprule
\textbf{Component} & \textbf{Power (W)} \\
\midrule
CPU & 16.7 \\
GPU & 22.9 \\
ANE & 0.0 \\
\midrule
\textbf{Total} & \textbf{39.6} \\
\bottomrule
\end{tabular}
\end{center}

This yields an efficiency of \textbf{0.51 tok/s per watt}. The ANE power draw of 0W independently confirms that the Neural Engine is completely unused during autoregressive decode, validating our analysis in Section~\ref{sec:nax} that NAX provides no benefit for M=1 inference.

\section{Discussion}

\subsection{The Shifting Bottleneck}

Our optimization journey reveals a classic systems phenomenon: eliminating one bottleneck exposes the next. The progression:

\begin{enumerate}
    \item \textbf{M3 Max baseline:} Expert I/O (47\%) $\rightarrow$ optimize I/O
    \item \textbf{Temporal prediction:} I/O hidden, GPU matmul (42\%) $\rightarrow$ optimize kernels
    \item \textbf{Fused dispatch:} GPU encoding (38\%) $\rightarrow$ pre-encode CMD2
    \item \textbf{Current:} Dense matmuls dominate, limited by memory bandwidth
\end{enumerate}

Notably, the theoretical I/O ceiling for 2-bit experts is 943~MB $\div$ 17.5~GB/s = 53.9ms = 18.6~tok/s. Our system achieves 20.34~tok/s---\textbf{surpassing the 2-bit I/O ceiling} through temporal prediction and fused scheduling that overlap I/O with compute.

The remaining bottleneck (GPU dense matmuls at 42\%) is memory-bandwidth-limited on unified memory architecture. Further improvement requires either:
\begin{itemize}
    \item Lower-precision dense weights (currently 4-bit)
    \item Hardware with higher memory bandwidth
    \item Algorithmic changes (sparse attention, pruning)
\end{itemize}

\subsection{Temporal Prediction: Why It Works}

The 27\% cross-token expert correlation surprised us---MoE routing is typically described as highly dynamic. We attribute this to:

\begin{enumerate}
    \item \textbf{Semantic coherence:} Adjacent tokens in a sentence share meaning
    \item \textbf{Attention locality:} Hidden states evolve smoothly
    \item \textbf{Expert specialization:} Some experts handle common patterns
\end{enumerate}

Importantly, the correlation is cross-token but not cross-layer. Each layer's routing depends on its specific input, which is transformed by intervening computations.

\subsection{GGUF Quantization: Precision Where It Matters}

Unsloth's selective precision (BF16 for layer 27, Q5\_K for outlier experts) demonstrates that uniform quantization leaves quality on the table. The open question is whether these outliers can be identified automatically or require per-model analysis.

\subsection{Autoresearch Methodology}

The 36-experiment trajectory (78\% discard rate) underscores the value of systematic exploration. Key practices:
\begin{itemize}
    \item Every experiment committed to git before execution
    \item Quality gate (PPL threshold) enforced automatically
    \item Immediate revert on regression
    \item Both positive and negative results logged
\end{itemize}

Human-AI collaboration was essential: the human provided systems intuition ("trust the OS cache"), while the AI implemented and benchmarked systematically.

\subsection{Limitations}

Our work has several limitations that should be considered when interpreting results:

\begin{enumerate}
    \item \textbf{Single model evaluation.} All 36 experiments were conducted exclusively on Qwen3.5-397B-A17B. Generalization to other MoE architectures---such as DeepSeek-V3 (671B, 37B active) or Mixtral (8$\times$7B)---is unverified. These models differ in expert count, activation ratio, and routing mechanism, any of which could affect the efficacy of temporal prediction or quantization strategies.

    \item \textbf{No standardized benchmark evaluation.} Model quality is measured solely via perplexity on WikiText-2 (499 and 2000 tokens). We have not evaluated on established benchmarks such as MMLU, GPQA, or SWE-bench. Perplexity correlates with downstream task performance but does not guarantee it, particularly after quantization.

    \item \textbf{Prefill performance.} Prefill throughput remains at 5.52 tok/s (prompt processing), making long-context use cases---such as document summarization or retrieval-augmented generation---impractical at interactive speeds. The prefill bottleneck stems from expert I/O: each prompt token requires its own set of expert reads, and temporal prediction cannot help (there is no prior-token routing history). ANE co-processing for batch prefill (Section~\ref{sec:nax}) is the most promising path forward.

    \item \textbf{Single hardware platform.} All measurements are on a single M5 Max (128GB) configuration. Results may not generalize to other Apple Silicon chips---M4 Pro/Max (no NAX, different SSD bandwidth), M3 Ultra (higher memory bandwidth but different GPU core count), or future M-series hardware. The 17.5 GB/s SSD bandwidth and 128GB unified memory are specific to this configuration and directly influence the effectiveness of our I/O optimizations.

    \item \textbf{Quality evaluation scope.} Output quality was evaluated via perplexity metrics only. Q3 quantization at this scale produces acceptable quality for short-context conversational tasks (under $\sim$200 tokens output) but exhibits degradation on long-form generation, complex reasoning, and extended lists. Standardized capability benchmarks were not evaluated. The primary contribution of this work is inference speed optimization, not model quality improvement.
\end{enumerate}

\section{Conclusion}

We demonstrate that Qwen3.5-397B-A17B runs at 20.34 tok/s on Apple M5 Max---a \textbf{4.67$\times$ improvement} over the prior state-of-the-art of 4.36 tok/s on M3 Max. The key innovations are:

\begin{enumerate}
    \item \textbf{Temporal expert prediction} exploits 27\% cross-token routing correlation to hide 65\% of expert I/O latency
    \item \textbf{GGUF Q3 quantization} reduces expert size by 19\% with importance-aware mixed precision
    \item \textbf{Fused command scheduling} eliminates CPU-GPU synchronization through pre-encoded command buffers
\end{enumerate}

We also establish that NAX provides no benefit for M=1 decode (tile padding overhead) and that extreme quantization (1-bit, ternary) catastrophically fails for post-training MoE compression.

The bottleneck has shifted from expert I/O (47\% $\rightarrow$ 10\%) to GPU dense matmuls (42\%), which are memory-bandwidth-limited. Future work should explore ANE co-processing for batch prefill and lower-precision dense weights.

\paragraph{Code Availability.} The complete Flash-MoE implementation, autoresearch harness, and experiment logs are available at \url{https://github.com/Anemll/flash-moe}.

\section*{Acknowledgments}

We thank the Unsloth team for their GGUF quantization work, the llama.cpp project for Metal kernel references, and the MLX team at Apple for their framework.

\bibliographystyle{plain}
\begin{thebibliography}{20}

\bibitem{flashmoe2025}
Daniel Woods.
\newblock Flash-MoE: Streaming a 397B Parameter Mixture-of-Experts Model from NVMe at 5.7 Tokens/Second on Consumer Hardware.
\newblock Technical report, 2025.

\bibitem{llminflash}
Keivan Alizadeh, Iman Mirzadeh, Dmitry Belenko, et al.
\newblock LLM in a Flash: Efficient Large Language Model Inference with Limited Memory.
\newblock \emph{arXiv preprint arXiv:2312.11514}, 2023.

\bibitem{qwen3}
An Yang, Baosong Yang, Beichen Zhang, et al.
\newblock Qwen3 Technical Report.
\newblock Technical report, Alibaba Group, 2025.

\bibitem{mixtral}
Albert Q. Jiang, Alexandre Sablayrolles, et al.
\newblock Mixtral of Experts.
\newblock \emph{arXiv preprint arXiv:2401.04088}, 2024.

\bibitem{deepseek}
DeepSeek-AI.
\newblock DeepSeek-V3 Technical Report.
\newblock Technical report, 2024.

\bibitem{deltanet}
Songlin Yang, Bailin Wang, Yikang Shen, et al.
\newblock Gated Delta Networks: Improving Mamba2 with Delta Rule.
\newblock \emph{arXiv preprint arXiv:2412.06464}, 2024.

\bibitem{gptq}
Elias Frantar, Saleh Ashkboos, Torsten Hoefler, and Dan Alistarh.
\newblock GPTQ: Accurate Post-Training Quantization for Generative Pre-trained Transformers.
\newblock \emph{arXiv preprint arXiv:2210.17323}, 2022.

\bibitem{awq}
Ji Lin, Jiaming Tang, Haotian Tang, et al.
\newblock AWQ: Activation-aware Weight Quantization for LLM Compression and Acceleration.
\newblock \emph{arXiv preprint arXiv:2306.00978}, 2023.

\bibitem{quip}
Jerry Chee, Yaohui Cai, Volodymyr Kuleshov, and Christopher De Sa.
\newblock QuIP: 2-Bit Quantization of Large Language Models With Guarantees.
\newblock \emph{arXiv preprint arXiv:2307.13304}, 2023.

\bibitem{bitnet}
Hongyu Wang, Shuming Ma, Li Dong, et al.
\newblock BitNet: Scaling 1-bit Transformers for Large Language Models.
\newblock \emph{arXiv preprint arXiv:2310.11453}, 2023.

\bibitem{mlx}
Awni Hannun, Jagrit Digani, et al.
\newblock MLX: Efficient and Flexible Machine Learning on Apple Silicon.
\newblock Technical report, Apple, 2023.

\bibitem{llamacpp}
Georgi Gerganov et al.
\newblock llama.cpp: Port of Facebook's LLaMA model in C/C++.
\newblock \url{https://github.com/ggerganov/llama.cpp}, 2023.

\bibitem{unsloth}
Daniel Han and Michael Han.
\newblock Unsloth: Fast and Memory-Efficient LLM Fine-tuning and Quantization.
\newblock \url{https://github.com/unslothai/unsloth}, 2024.

\bibitem{turboquant}
Amir Zandieh, Insu Han, Majid Daliri, and Amin Karbasi.
\newblock TurboQuant: Online Vector Quantization with Sub-quadratic Encoding.
\newblock \emph{arXiv preprint arXiv:2504.19874}, 2025.

\bibitem{specdecoding}
Yaniv Leviathan, Matan Kalman, and Yossi Matias.
\newblock Fast Inference from Transformers via Speculative Decoding.
\newblock \emph{arXiv preprint arXiv:2211.17192}, 2022.

\bibitem{dynamicmoe}
Ping Huang, Xin Wang, et al.
\newblock Harder Tasks Need More Experts: Dynamic Routing in MoE Models.
\newblock \emph{arXiv preprint arXiv:2403.07652}, 2024.

\bibitem{expertchoice}
Yanqi Zhou, Tao Lei, Hanxiao Liu, et al.
\newblock Mixture-of-Experts with Expert Choice Routing.
\newblock \emph{arXiv preprint arXiv:2202.09368}, 2022.

\end{thebibliography}

\appendix

\section{Complete Experiment Log}

Table~\ref{tab:experiments} presents the full autoresearch trajectory.

\begin{table*}[h]
\centering
\caption{Complete autoresearch experiment log (M5 Max). Status: keep = improvement, discard = regression or failure. Rows numbered sequentially 1--36 matching \texttt{results.tsv}.}
\label{tab:experiments}
\footnotesize
\begin{tabular}{rllrrrp{6cm}}
\toprule
\textbf{\#} & \textbf{Commit} & \textbf{Status} & \textbf{tok/s} & \textbf{$\Delta$\%} & \textbf{PPL} & \textbf{Description} \\
\midrule
1 & 26cd7f8 & keep & 10.61 & +15.2 & 5.62 & Starting point: 4bit + Q8\_0 embed + Q6\_K LM head \\
2 & ee0d861 & keep & 10.66 & +0.5 & 5.62 & cache-io-split=4 default \\
3 & a0fa67c & discard & 10.26 & -3.3 & 5.58 & Q3 auto-detect (score regression) \\
4 & e06a0fa & discard & 9.57 & -9.8 & $>$1M & NAX Q6\_K prequant (430ms startup) \\
5 & a145103 & keep & 12.11 & +14.1 & 5.62 & 16 IO threads (max parallelism) \\
6 & e212e90 & discard & 11.60 & +9.3 & 5.62 & split=8 + 32 threads (syscall overhead) \\
7 & fdae1de & discard & 12.00 & +13.1 & 5.62 & v3\_8k kernel (occupancy limited) \\
8 & 95f4f72 & discard & 11.22 & +5.8 & 5.62 & Single-encoder batch (no speculation) \\
9 & 7fb181e & keep & 16.40 & +54.6 & 5.62 & Temporal prediction default-on \\
10 & b2e608e & discard & 12.06 & +13.7 & 5.62 & 2-token union (SSD bandwidth overflow) \\
11 & ea7d038 & discard & 14.27 & +34.5 & 5.62 & v7k kernel (50\% occupancy) \\
12 & stashed & discard & 14.23 & +34.1 & 5.62 & Pred-path reorder (compiler regression) \\
13 & stashed & discard & 10.84 & +2.1 & 5.62 & GPU gate in CMD1 (delta-net bug) \\
14 & stashed & discard & 9.37 & -11.7 & 5.62 & 4-token history (21.6\% hit rate) \\
15 & stashed & discard & 8.65 & -18.5 & 5.62 & Miss-read split (syscall overhead) \\
16 & 3ddeb82 & keep & 18.67 & +76.0 & 5.58 & Q3 experts + temporal \\
17 & stashed & discard & 16.25 & +53.2 & 5.58 & Fused residual+norm (occupancy) \\
18 & stashed & discard & 18.59 & +75.2 & 5.58 & 32 threads/TG (no improvement) \\
19 & stashed & discard & 14.30 & +34.8 & 5.58 & Cross-layer prediction (0\% hit rate) \\
20 & stashed & discard & 17.00 & +60.2 & 5.58 & FMA on fast kernel (register pressure) \\
21 & stashed & discard & 18.14 & +71.0 & 5.62 & Dequant-first FMA (serial dependency) \\
22 & stashed & discard & 13.43 & +26.6 & 8.09 & 2-bit + temporal (PPL fails) \\
23 & stashed & discard & 14.28 & +34.6 & 5.59 & Half x\_shared (conversion overhead) \\
24 & stashed & discard & 16.56 & +56.1 & 5.58 & split=1 with Q3 (less parallelism) \\
25 & stashed & discard & 18.54 & -0.7 & 5.58 & CMD1+CMD3 fusion (no improvement) \\
26 & stashed & discard & 18.39 & -1.5 & 5.63 & 4-bit MLX LM head (cache instability) \\
27 & 0a3b923 & keep & 19.11 & +80.1 & 5.58 & CMD2 pre-encode GDN \\
28 & stashed & discard & 13.74 & +29.5 & 5647 & 1-bit QJL experts (catastrophic) \\
29 & stashed & discard & 18.07 & +70.3 & 11.49 & Ternary 2-bit (84\% sparsity) \\
30 & stashed & discard & 17.41 & +64.1 & 6.54 & K=3 routing (quality collapse) \\
31 & stashed & discard & 16.30 & -14.7 & 5.58 & K+1 prediction (bandwidth overflow) \\
32 & stashed & discard & 17.58 & +65.7 & 5.58 & Miss-read overlap (deadlock) \\
33 & stashed & discard & 19.18 & +70.0 & 5.58 & Q6\_K LM head v2 (barrier overhead) \\
34 & 2dc5b66 & keep & 19.87 & +87.3 & 5.58 & Fused Q/K/V projection dispatch \\
35 & 9c7ed1a & keep & 20.34 & +91.7 & 5.58 & CMD2 pre-encode for full-attention \\
36 & stashed & discard & 19.51 & +84.0 & 5.58 & CMD1 pre-encode (driver overhead) \\
\bottomrule
\end{tabular}
\end{table*}

\section{GGUF Block Layouts}

The GGUF quantization formats used:

\begin{lstlisting}[language=C,caption=GGUF block structures]
// Q8_0: 32 weights per block
struct block_q8_0 {
    half d;           // 2 bytes
    int8_t qs[32];    // 32 bytes
};  // 34 bytes / 32 weights

// Q6_K: 256 weights per block
struct block_q6_k {
    uint8_t ql[128];  // lower 4 bits
    uint8_t qh[64];   // upper 2 bits
    int8_t scales[16];
    half d;
};  // 210 bytes / 256 weights

// IQ3_XXS: 256 weights per block
struct block_iq3_xxs {
    half d;
    uint8_t qs[3*32]; // 2-bit grid indices
    uint8_t signs[32];
};  // 130 bytes / 256 weights

// IQ4_XS: 256 weights per block
struct block_iq4_xs {
    half d;
    uint16_t scales_h;
    uint8_t scales_l[8];
    uint8_t qs[128];
};  // 148 bytes / 256 weights
\end{lstlisting}

\clearpage
\section{Hardware Specifications}

\begin{table}[h]
\centering
\caption{Hardware comparison: M3 Max vs M5 Max}
\label{tab:hardware}
\begin{tabular}{lrr}
\toprule
\textbf{Specification} & \textbf{M3 Max} & \textbf{M5 Max} \\
\midrule
CPU cores & 16 (12P+4E) & 16 (12P+4E) \\
GPU cores & 40 & 40 \\
Memory & 48 GB & 128 GB \\
Memory BW & 400 GB/s & 500 GB/s \\
SSD BW & 17.5 GB/s & 17.5 GB/s \\
Metal version & 3.2 & 4.0 \\
NAX & No & Yes \\
Inference power & --- & 39.6 W \\
Efficiency & --- & 0.51 tok/s/W \\
\bottomrule
\end{tabular}
\end{table}

\end{document}
